\section{Discussion}\label{sec:discussion}

\subsection{What the joint view shows that single-axis evaluations miss}
Evaluated one axis at a time, our models look deployable. In-domain AUROC is near 1, calibration error is near 0, and the thresholds hit their nominal false-positive rate. Viewed jointly, three problems appear that a single-axis evaluation would not reveal.
\begin{itemize}
\item \emph{Calibration does not transfer.} It is a property of the pair of source and target distributions, not of the model. Post-hoc recalibration on source data, the usual fix, leaves cross-dataset ECE at 0.28--0.45. Well-calibrated in-domain probabilities should therefore not be used to plan alert volumes on a new network, consistent with observations outside security~\cite{ovadia2019trust}.
\item \emph{Ranking metrics hide operational failure.} The same average AUROC can come with a thirtyfold alarm flood (CSE-CIC-IDS2018$\to$CIC-IDS2017) or with near-total silence (ACI-IoT-2023$\to$CIC-IDS2017). This extends Jin and Lee's observation of threshold collapse~\cite{jin2026benchmark} to a failure mode that is arguably worse, because a silent detector raises no suspicion.
\item \emph{Unseen-family misses are a thresholding problem more than a representation problem.} The held-out family usually ranks above benign traffic but below the operating threshold. Almost a third of families are missed, and almost a quarter (11 of 47) are missed with high calibrated confidence, so uncertainty-based rejection alone~\cite{talpini2024enhancing} would not flag them.
\end{itemize}

\subsection{Implications for practice}
Four recommendations follow from the results.
\begin{enumerate}
\item \emph{Fix volume, not score.} Define the alert policy as a review budget and enforce it with a quantile of the priority on recent unlabeled traffic. This costs nothing and kept the median queue size within 5\% of plan under every shift we tested, whereas thresholds carried over from training were off by factors of 0.4 to 54.
\item \emph{Treat novelty evidence as insurance.} Fusing a benign-referenced novelty p-value with the detector p-value (Fisher's method) roughly doubles recall at a 1\% budget on the families the detector misses, at a cost of 2--6 points on families it already detects. Whether that trade is worth it depends on the relative cost of a missed campaign and a slightly less efficient queue; in many SOCs a missed zero-day is the costlier error. A practical middle ground is to reserve a small fixed share of the budget for novelty-ranked flows. This is the budget-splitting analogue of the Bonferroni rule and a natural next step.
\item \emph{Do not expect any post-hoc layer to rescue cross-network transfer.} Neither calibration nor novelty fusion recovered cross-dataset recall; the models had learned features specific to how each capture was generated. Local retraining, or at least local benign reference data for the novelty score, remains necessary.
\item \emph{Report at realistic base rates.} Precision at 1\% prevalence was at most 0.49 even in-domain, and about 0.03 across datasets. Balanced-test accuracy says little about analyst workload~\cite{axelsson2000base,alahmadi2022falsepositives}.
\end{enumerate}

\subsection{Choice of novelty reference}
The novelty score should measure distance to \emph{benign} traffic, as the null hypothesis requires. In a first complete run of Protocol OS we used all training flows, including attacks, as the kNN reference. Because unseen attacks often resemble known ones, that choice made them look typical: the mean AUROC between held-out-family flows and benign flows was 0.74, against 0.83 with the benign reference (both runs are released). The difference was largest on UNSW-NB15, where the median rose from 0.52 to 0.91. We report the benign-reference variant throughout. The kNN score is also blind to some attacks. Floods and sweeps in ACI-IoT-2023 consist of very short, highly regular flows that lie \emph{closer} to benign traffic than benign flows lie to each other (novelty AUROC below 0.5 for SYN flood and ping sweep in a pilot). Detecting such hyper-typical traffic needs volumetric or host-level context that per-flow features do not carry.

\subsection{Threats to validity}\label{sec:threats}
\emph{Data availability.} Network restrictions in our execution environment prevented us from obtaining the NetFlow-v2 versions of several datasets~\cite{sarhan2022standard}, the full CIC-IDS2017 release and the full ToN-IoT normal class. We used a published subsample of CIC-IDS2017 and a ToN-IoT copy with 50{,}000 of the 300{,}000 normal records; the ACI-IoT-2023 file came from a public mirror whose class counts match the official release. All provenance, including checksums, is released. Because all operational metrics are computed on streams resampled to a fixed prevalence, these choices affect the diversity of benign training data but not the definitions of the metrics. Known labelling problems in CIC-IDS2017 and CSE-CIC-IDS2018~\cite{engelen2021troubleshooting,lanvin2023errors,flood2024bad} also affect our results; we did not relabel the data.

\emph{Feature alignment.} The three CICFlowMeter datasets were produced with different tool versions, and some features differ in implementation, not only in traffic. This is part of what a practitioner faces when moving a model between networks, but it means our cross-dataset results mix network shift with tooling shift.

\emph{Scope of models and splits.} We used four standard tabular model families, which remain strong on this kind of data~\cite{grinsztajn2022tree}, and random stratified splits within datasets. Temporal splits~\cite{pendlebury2019tesseract,luay2025temporal}, deep sequence or graph models, and adaptive conformal methods for drifting streams~\cite{gibbs2021adaptive,barber2023conformal} are natural extensions. The unseen-family protocol used LightGBM only, and the grouping in Table~\ref{tab:gain} is post hoc.

\emph{Guarantees.} The conformal false-alarm guarantee holds only under exchangeability of benign target and calibration flows. Our cross-dataset results show that this assumption fails across networks, which is why we fix the budget empirically rather than relying on the nominal $\alpha$.
